\clearpage
\VerdierPhysicalPageAnchor{0229}%
\begingroup\small\setlength{\parskip}{0pt}\setlength{\emergencystretch}{1.5em}
\setlength{\abovedisplayskip}{4pt plus 2pt minus 2pt}%
\setlength{\belowdisplayskip}{4pt plus 2pt minus 2pt}%

\VerdierSemanticAnchor{ch4d-4-1-6-cont}
\noindent 如下定义的 $n$-重复形：

\VerdierSemanticAnchor{ch4d-4-1-6-1}
\begin{equation}
({}^{m}\!X(p,q))^\sigma=
\begin{cases}
X^\sigma,&-q\leq\sigma_m<-p,\\
0,&\sigma_m<-q\ \lor\ -p\leq\sigma_m,
\end{cases},
\tag{4.1.6.1}
\end{equation}

\noindent ${}^{m}\!X(p,q)$ 的微分就是 $X$ 的微分，除非所考察微分的源或靶
是零对象，在这种情形可能例外。

对 $\mathcal{F}\ell_1(\widetilde{\mathbb Z})$ 中满足 $p\leq p'$ 且
$q\leq q'$ 的任意两个元素 $(p,q)$ 与 $(p',q')$，以

\VerdierSemanticAnchor{ch4d-4-1-6-2}
\begin{equation}
f\bigl((p,q),(p',q')\bigr)\colon
{}^{m}\!X(p,q)\longrightarrow{}^{m}\!X(p',q')
\tag{4.1.6.2}
\end{equation}

表示如下 $n$-重复形态射：其分量是恒等态射，除非源或靶是零对象，在这种
情形可能例外。最后，对
$\mathcal{F}\ell_2(\widetilde{\mathbb Z})$ 的每个元素 $(p,q,r)$，以

\VerdierSemanticAnchor{ch4d-4-1-6-3}
\begin{equation}
g(p,q,r)\colon{}^{m}\!X(q,r)
\longrightarrow T_m\bigl({}^{m}\!X(p,q)\bigr)
\tag{4.1.6.3}
\end{equation}

表示如下态射（$T_m$ 表示沿方向 $m$ 的平移函子（第一章，2.3.1）），其分量为：

\VerdierSemanticAnchor{ch4d-4-1-6-4}
\begin{equation}
\left\{
\begin{aligned}
g(p,q,r)^\sigma&=0,
  &&\sigma_m\neq-q-1,\\
g(p,q,r)^\sigma&=-d_X^{\sigma,m},
  &&\sigma_m=-q-1.
\end{aligned}
\right.
\tag{4.1.6.4}
\end{equation}

现在引入新记号。首先，滥用记号，以
$\int_n\colon\operatorname{comp}^n(\mathcal A)\longrightarrow
\operatorname{comp}(\mathcal A)$ 表示关联单复形函子的 $\Pi$-扩张或
$\Sigma$-扩张（前提是 $\mathcal A$ 中存在可数积或可数和（第一章，
2.2.3））。置

\VerdierSemanticAnchor{ch4d-4-1-6-5}
\begin{equation}
\left\{
\begin{aligned}
\int_n{}^{m}\!X(p,q)
  &={}^{m}\!Y(p,q),\\
\int_n{}^{m}\!X(-\infty,q)
  &={}^{m}\!Y(-\infty,q)={}^{m}\!Y(q),\\
\int_n{}^{m}\!X(-\infty,+\infty)=\int_n X
  &={}^{m}\!Y(-\infty,+\infty)={}^{m}\!Y(+\infty)=Y.
\end{aligned}
\right.
\tag{4.1.6.5}
\end{equation}
\endgroup
